\section{精确闭合：四个等价判据与最小任务商}
\label{app:exact-theory}

本附录证明有限确定性粗粒化的精确理论。核心量词是：同一宏观初态必须对
\emph{所有}微观初态分布给出同一宏观下一步分布。这正是本文使用 strong
lumpability、而非只对某一类初始分布成立的弱版本的原因。

\subsection{因子分解的四个等价形式}

设 $q_X:X\twoheadrightarrow\bar X$、$q_Y:Y\twoheadrightarrow\bar Y$ 为
满射，$K\in\Stoch(X,Y)$。

\begin{theorem}[有限确定性粗粒化的四个等价判据]
\label{thm:app-four-equivalences}
以下命题等价。
\begin{enumerate}
  \item 存在 $\bar K\in\Stoch(\bar X,\bar Y)$ 使
  \begin{equation}
    KQ_{q_Y}=Q_{q_X}\bar K.
    \label{eq:app-exact-square}
  \end{equation}
  该 $\bar K$ 必唯一。
  \item 若 $q_X(x)=q_X(x')$，则对每个 $\bar y\in\bar Y$，
  \begin{equation}
    K\bigl(x,q_Y^{-1}(\bar y)\bigr)
    =K\bigl(x',q_Y^{-1}(\bar y)\bigr).
    \label{eq:app-block-probability}
  \end{equation}
  等价地，上式对任意 $\bar A\subseteq\bar Y$ 的原像成立。
  \item 粗粒化观测空间在 $P_K$ 下闭合：
  \begin{equation}
    P_K\bigl(q_Y^*\mathbb R^{\bar Y}\bigr)
    \subseteq q_X^*\mathbb R^{\bar X}.
    \label{eq:app-observable-invariance}
  \end{equation}
  \item 对任意 $\mu,\nu\in\Delta(X)$，
  \begin{equation}
   q_{X*}\mu=q_{X*}\nu
   \quad\Longrightarrow\quad
   q_{Y*}(\mu K)=q_{Y*}(\nu K).
   \label{eq:app-distribution-autonomy}
  \end{equation}
\end{enumerate}
\end{theorem}

\begin{proof}
$1\Rightarrow2$：由式~\eqref{eq:app-exact-square}，
\[
 K\bigl(x,q_Y^{-1}(\bar y)\bigr)
 =(KQ_{q_Y})(x,\bar y)
 =\bar K(q_X(x),\bar y),
\]
故该值只依赖 $q_X(x)$。

$2\Rightarrow1$：对 $\bar x\in\bar X$ 任选 $x\in q_X^{-1}(\bar x)$，定义
\[
  \bar K(\bar x,\bar y)
  :=K\bigl(x,q_Y^{-1}(\bar y)\bigr).
\]
条件~\eqref{eq:app-block-probability} 保证定义与代表元无关；非负性和行和为
$1$ 来自 $K$ 与 $q_Y$ 的满射性。逐行即得
式~\eqref{eq:app-exact-square}。若另有 $\bar K'$ 满足该式，则对任意
$\bar x$ 选取 $q_X(x)=\bar x$，有
$\bar K(\bar x,\cdot)=(KQ_{q_Y})(x,\cdot)=\bar K'(\bar x,\cdot)$，
故唯一。

$1\Rightarrow3$：对 $\bar f\in\mathbb R^{\bar Y}$，
\[
 P_K(q_Y^*\bar f)
 =P_{KQ_{q_Y}}\bar f
 =P_{Q_{q_X}\bar K}\bar f
 =q_X^*(P_{\bar K}\bar f).
\]

$3\Rightarrow2$：取 $\bar f=\one_{\{\bar y\}}$。则
\[
  P_K(q_Y^*\bar f)(x)
  =K\bigl(x,q_Y^{-1}(\bar y)\bigr)
\]
按条件在 $q_X$ 的纤维上常值。对多个 $\bar y$ 求和即得到事件版本。

$1\Rightarrow4$：对任意 $\mu$，
\[
  q_{Y*}(\mu K)=\mu KQ_{q_Y}
  =\mu Q_{q_X}\bar K=(q_{X*}\mu)\bar K.
\]
因此相同的宏观初始分布产生相同的宏观输出分布。

$4\Rightarrow2$：对同一 $q_X$ 纤维中的 $x,x'$，取
$\mu=\delta_x$、$\nu=\delta_{x'}$，再读取每个 $\bar y$ 的概率即可。
\end{proof}

\begin{corollary}[齐次系统的全过程投影]
\label{cor:app-exact-process}
若 $K\in\Stoch(X,X)$、$q:X\twoheadrightarrow\bar X$ 且
$KQ_q=Q_q\bar K$，则
\[
  K^nQ_q=Q_q\bar K^n\qquad(n\ge0).
\]
对任意微观初始分布，以 $K$ 演化的链 $(X_t)$ 的投影 $(q(X_t))$ 是以
$\bar K$ 为转移核的 Markov 链。
\end{corollary}

\begin{proof}
第一式由归纳直接得到。给定 $X_t=x$ 时，$q(X_{t+1})$ 的条件分布为
$\bar K(q(x),\cdot)$；在给定任意宏观历史且当前宏观状态为 $\bar x$ 后，
对纤维 $q^{-1}(\bar x)$ 上的微观条件分布取混合，结果仍为
$\bar K(\bar x,\cdot)$。这就是 Markov 性。
\end{proof}

\subsection{有限观测代数就是分区代数}

称 $\cA\subseteq\mathbb R^X$ 为观测代数，如果它是包含常数函数且对点乘
封闭的线性子空间。定义
\[
 x\sim_{\cA}x'
 \quad\Longleftrightarrow\quad
 f(x)=f(x')\quad\forall f\in\cA,
\]
并记商映射为 $q_{\cA}:X\twoheadrightarrow X_{\cA}$。

\begin{lemma}[有限观测代数--分区对应]
\label{lem:app-finite-algebra-partition}
对有限 $X$，
\begin{equation}
  \cA=q_{\cA}^*\mathbb R^{X_{\cA}},
  \qquad
  \dim\cA=|X_{\cA}|.
  \label{eq:app-algebra-partition}
\end{equation}
反之，每个分区商 $q:X\twoheadrightarrow Y$ 都给出观测代数
$q^*\mathbb R^Y$。两个满射给出同一代数，当且仅当其纤维分区相同。
\end{lemma}

\begin{proof}
按定义，$\cA$ 中每个函数都在 $\sim_{\cA}$ 的等价类上常值，故
$\cA\subseteq q_{\cA}^*\mathbb R^{X_{\cA}}$。反过来，固定一个等价类
$C$。对每个其他等价类 $D$，可取 $f_D\in\cA$ 使其在 $C,D$ 上的常值
$a_D,b_D$ 不同。仿射多项式
\[
  h_D:=\frac{f_D-b_D\one}{a_D-b_D}\in\cA
\]
在 $C$ 上为 $1$、在 $D$ 上为 $0$。有限乘积
$\prod_{D\ne C}h_D$ 就是 $\one_C$；只有一个等价类时取空乘积 $\one$。
因此所有块指示函数均在 $\cA$ 中，任何块上常值函数也在其中，得到第一式。
拉回沿满射为单射，块指示函数又线性无关，故维数等式成立。反向陈述直接由
点乘和常数闭合得到。
\end{proof}

\begin{theorem}[闭合观测代数与闭合商]
\label{thm:app-closed-algebra}
设 $K\in\Stoch(X,X)$，$\cA\subseteq\mathbb R^X$ 是观测代数。则
\[
  P_K\cA\subseteq\cA
\]
当且仅当存在唯一 $K_{\cA}\in\Stoch(X_{\cA},X_{\cA})$ 使
\[
  KQ_{q_{\cA}}=Q_{q_{\cA}}K_{\cA}.
\]
\end{theorem}

\begin{proof}
由引理~\ref{lem:app-finite-algebra-partition}，
$\cA=q_{\cA}^*\mathbb R^{X_{\cA}}$。所述不变性正是
定理~\ref{thm:app-four-equivalences} 的条件~3。
\end{proof}

常数代数和 $\mathbb R^X$ 对每个 $P_K$ 都闭合，分别对应一状态商与恒等商。
所以``某个闭合商存在''本身没有非平凡内容；必须指定不可丢弃的任务观测。

\subsection{核族的任务相对最小精确商}

正文要求同一个表示服务于一族允许动力学，而允许每个微观 kernel 有自己的
宏观 kernel。设 $g:X\twoheadrightarrow O$ 是非恒定任务观测，
$\cK\subseteq\Stoch(X,X)$ 非空。令
\begin{equation}
  \cA_0:=g^*\mathbb R^O,
  \qquad
  \cA_{n+1}:=
  \operatorname{alg}\!\left(
    \cA_n\cup\bigcup_{K\in\cK}P_K\cA_n
  \right).
  \label{eq:app-task-algebra-iteration}
\end{equation}

\begin{theorem}[核族的任务相对最小精确宏观模型]
\label{thm:app-minimal-task-quotient}
迭代~\eqref{eq:app-task-algebra-iteration} 在至多 $|X|-|O|$ 次严格增加后
稳定；记稳定代数为 $\cA_*$，相应商为
$q_*:X\twoheadrightarrow X_*$。则：
\begin{enumerate}
  \item $\cA_*$ 是包含 $\cA_0$、并对全部 $P_K$（$K\in\cK$）不变的
  最小观测代数；
  \item 存在唯一 $r_*:X_*\to O$，并且对每个 $K\in\cK$ 存在唯一
  $K_*^{(K)}\in\Stoch(X_*,X_*)$，使
  \begin{equation}
    KQ_{q_*}=Q_{q_*}K_*^{(K)},
    \qquad
    g=r_*\circ q_*;
    \label{eq:app-minimal-task-square}
  \end{equation}
  \item 若另有满射 $q:X\twoheadrightarrow Y$、读出 $r:Y\to O$，并且
  对每个 $K\in\cK$ 都有 $L_K\in\Stoch(Y,Y)$ 满足
  \[
    g=r\circ q,
    \qquad
    KQ_q=Q_qL_K,
  \]
  则存在唯一满射 $s:Y\twoheadrightarrow X_*$ 使
  \begin{equation}
    q_*=s\circ q,
    \qquad r=r_*\circ s,
    \qquad L_KQ_s=Q_sK_*^{(K)}\quad(K\in\cK).
    \label{eq:app-minimal-universal-property}
  \end{equation}
\end{enumerate}
换言之，$q_*$ 是全部任务保真、并对 $\cK$ 共同精确闭合的确定性商中
唯一最粗者。
\end{theorem}

\begin{proof}
$\cA_n$ 单调增加；每次严格增加至少使维数增加 $1$。由于
$\dim\cA_0=|O|$、$\dim\mathbb R^X=|X|$，至多发生
$|X|-|O|$ 次严格增加。稳定时
$P_K\cA_*\subseteq\cA_*$ 对每个 $K\in\cK$ 成立。若 $\mathcal B$ 是
任意包含 $\cA_0$ 且对全部 $P_K$ 不变的观测代数，则从
$\cA_0\subseteq\mathcal B$ 出发归纳得到
$\cA_n\subseteq\mathcal B$，故 $\cA_*\subseteq\mathcal B$。这证明
第一项。

逐个 $K$ 应用定理~\ref{thm:app-closed-algebra}，得到唯一
$K_*^{(K)}$。又因 $\cA_0\subseteq\cA_*$，$g$ 在 $q_*$ 的纤维上常值，
故定义 $r_*(q_*(x)):=g(x)$；它良定义且因 $q_*$ 满射而唯一。

对第三项，候选 $q$ 的分区代数
$\cA_q:=q^*\mathbb R^Y$ 包含 $\cA_0$，且由各个精确交换式和
定理~\ref{thm:app-four-equivalences} 对全部 $P_K$ 不变。最小性给出
$\cA_*\subseteq\cA_q$，所以存在唯一满射 $s$ 使 $q_*=s\circ q$；任务
读出的等式由 $q$ 满射立即得到。最后，对每个 $K\in\cK$，
\[
 Q_qL_KQ_s=KQ_qQ_s=KQ_{q_*}
 =Q_{q_*}K_*^{(K)}=Q_qQ_sK_*^{(K)}.
\]
满射 $q$ 的各行覆盖 $Y$ 的全部 Dirac 分布，可以消去左侧的 $Q_q$，得到
$L_KQ_s=Q_sK_*^{(K)}$。
\end{proof}

\begin{definition}[核族的精确任务相对闭合复杂度]
\label{def:app-exact-complexity}
定义
\begin{equation}
 c_{\cK}^{0}(g):=|X_*|=\dim\cA_*.
 \label{eq:app-exact-complexity}
\end{equation}
当 $\cK=\{K\}$ 时简记为 $c_K^0(g)$。
\end{definition}

定理~\ref{thm:app-minimal-task-quotient} 立即给出
\begin{equation}
 c_{\cK}^{0}(g)=
 \min\left\{|Y|:
 \begin{array}{l}
 \exists q:X\twoheadrightarrow Y,\ \exists r:Y\to O,\quad g=r\circ q,\\
 \forall K\in\cK\ \exists L_K\in\Stoch(Y,Y),\quad KQ_q=Q_qL_K
 \end{array}\right\}.
 \label{eq:app-exact-complexity-min}
\end{equation}
因此，任务相对的严格共同精确压缩存在，当且仅当
\begin{equation}
  |O|\le c_{\cK}^{0}(g)<|X|.
  \label{eq:app-nontrivial-exact}
\end{equation}
若按本文约定 $g$ 非恒定，则左端至少为 $2$。该结论不保证右侧严格不等式：
一般有限 Markov 核族可能迫使任务分区一直细化到离散分区。

迭代的分区版本从 $g$ 的纤维开始，并反复按
\[
 x\equiv_{n+1}x'
 \Longleftrightarrow
 \left\{
 \begin{array}{l}
 x\equiv_nx',\\
 K(x,C)=K(x',C)
 \quad\forall K\in\cK,\ \forall C\in\Pi_n
 \end{array}\right.
\]
细化；命题~\ref{prop:app-stable-refinement} 的同一证明表明它有限终止于
$q_*$。这给出有限输入上的构造算法，但没有给出大状态空间中的计算效率保证。
